\chapter{Metals}\label{ch:m3:metals}

In the early hours of Tuesday 8 March 2022 the three-month nickel contract on the London
Metal Exchange rose from slightly below 50{,}000 dollars a tonne to 101{,}365 dollars at
06:08, more than doubling in about five hours; the day before it had opened just under
30{,}000. At 08:15 the exchange suspended trading. At 12:05 it cancelled every nickel trade
done since midnight. Its clearing house had calculated that, had the trades stood, members
would have had to post some 19.75 billion dollars of margin at short notice, with a risk of
multiple defaults. Base metals trade on an exchange whose contracts, warehouses and rules
differ from every other futures market, and the week of the nickel squeeze tested all of
them. This chapter covers the London metal market's \omterm{def:m3:metals:prompt}{prompt dates} and prices, its warehouse
system, the nickel episode, gold's two markets in London and New York, and iron ore.

\section{The London metal market: prompt dates and prices}

The London Metal Exchange trades aluminium, copper, zinc, nickel, lead, tin and other
metals. Its contracts are not monthly futures: they are forwards for delivery on any of a
dense set of dates, a structure inherited from the time it took ships to bring metal to
London.

\begin{definition}[Prompt date, three-month price]\label{def:m3:metals:prompt}
A \emph{prompt date}\index{prompt date} is a date on which an LME contract can be settled by
delivery of warrants against payment. There is one for every business day from tomorrow
(tom) to three months ahead, one every Wednesday from there to six months, and one on the
third Wednesday of each month beyond, up to a limit that depends on the metal. The
\emph{three-month price}\index{three-month price} (3M) is the price for the prompt date about
three months from the trade date; it is the most liquid LME contract and the reference for
most hedging.
\end{definition}

\begin{omfigure}
\begin{tikzpicture}
\begin{axis}[width=11cm,height=3.6cm,
  xlabel={months after the trade date},
  tick label style={font=\small},label style={font=\small},
  xmin=0,xmax=27,ymin=0.5,ymax=3.5,ytick={1,2,3},yticklabels={monthly,weekly,daily},
  grid=major,grid style={gray!25},
  table/col sep=comma]
\addplot[omDef,only marks,mark=|,mark size=3pt] table[x=months,y=row]{figdata/markets-3/08-metals/prompts.csv};
\end{axis}
\end{tikzpicture}

\omcaption{\omterm{def:m3:metals:prompt}{Prompt dates} available on 24 September 2026 under the structure of
\cref{def:m3:metals:prompt}, out to 27 months: 65 daily, 14 weekly and 21 monthly dates,
weekends only (no holiday calendar). Data: the chapter's build.}\label{fig:m3:metals:prompts}
\end{omfigure}

A dense date structure lets a consumer match its hedge to the day it needs metal, but it
spreads liquidity. Traders move positions between dates with carries, the metal market's
name for a calendar spread (One Quant Book 1, chapter 19), and the shortest of them is the
tom-next of One Quant Book 2, chapter 16: selling for tomorrow and buying back for the next
day lends metal for a day.

\begin{example}[A carry as an interest rate]\label{ex:m3:metals:carry}
Copper cash is at 9{,}800 dollars a tonne and three months at 9{,}880: an 80-dollar contango, 0.82\% of
the cash price for about 91 days, or 3.3\% a year. A holder of metal who sells cash and buys three
months (lends metal) gives up 80 dollars a tonne and saves the cost of financing and storing the
metal for three months; if money and storage cost more than 3.3\% a year, lending is the better
use of the metal. A backwardation of the same size would pay the lender 80 dollars instead: the
market pays for metal now.
\end{example}

\begin{definition}[LME official price]\label{def:m3:metals:official}
The \emph{LME official price}\index{LME official price} of a metal is the price set each
business day from the final bids and offers of the second session of the exchange's
open-outcry trading floor (the Ring), for the cash, three-month and some later dates; physical contracts across the industry price against
it, much as oil cargoes price against an agency's assessment.
\end{definition}

\section{Warehouses, warrants and queues}

LME contracts are delivered by transferring ownership of metal stored in warehouses the
exchange has approved, in documents that circulate like title.

\begin{definition}[Warehouse warrant, cancelled warrant, load-out queue]\label{def:m3:metals:warrant}
A \emph{warehouse warrant}\index{warehouse warrant} is a document of title to a specific lot of
metal of a specified brand and quality in an LME-approved warehouse, and the means of delivery
against LME contracts. A \emph{cancelled warrant}\index{cancelled warrant} is one whose owner has
told the warehouse that it will withdraw the metal: the lot leaves the LME system and waits to
be loaded out. The \emph{load-out queue}\index{load-out queue} is the backlog of cancelled metal
waiting to leave a warehouse when cancellations exceed its daily load-out capacity.
\end{definition}

\begin{omfigure}
\centering
\begin{tikzpicture}[font=\small,
  st/.style={draw,thick,rounded corners=2pt,align=center,minimum height=1cm,text width=2.4cm},
  arr/.style={-{Stealth[length=2.2mm]},thick}]
\node[st,draw=omExo,fill=omExo!8] (in) at (0,0) {metal delivered\\into a warehouse};
\node[st,draw=omDef,fill=omDef!8] (on) at (3.4,0) {on warrant:\\deliverable};
\node[st,draw=omProp,fill=omProp!8] (can) at (6.8,0) {cancelled\\warrant};
\node[st,draw=omThm,fill=omThm!8] (q) at (10.2,0) {load-out\\queue};
\draw[arr] (in) -- (on);
\draw[arr] (on) -- (can);
\draw[arr] (can) -- (q);
\draw[arr] (q.south) -- ++(0,-0.7) node[below,font=\footnotesize] {loaded out to a consumer};
\draw[arr] (on.south) -- ++(0,-0.7) node[below,font=\footnotesize] {traded on the LME};
\end{tikzpicture}

\omcaption{The life of a lot of metal in the LME system. Only metal on warrant can be
delivered against a contract; cancelled metal waits to leave, and a long queue ties up metal
that the market counts as stock. Schematic.}\label{fig:m3:metals:warrants}
\end{omfigure}

Stocks on warrant are what shorts can deliver; cancelled stocks are metal already claimed. A
fall in stocks on warrant, or a jump in cancellations, tightens the near dates and can turn the
curve into backwardation. Long queues have also been used to raise rent income and the premiums
consumers pay for prompt metal; the exchange's rules now cap the rent warehouses may charge on
metal in long queues.

\begin{definition}[LME lending rule]\label{def:m3:metals:lending}
An \emph{LME lending rule}\index{LME lending rule} requires a holder of a dominant long position
(of warrants and near-dated contracts) to lend (sell for the nearer date and buy back for the
later one) at no more than a set backwardation, so that holding most of the deliverable metal
cannot be used to squeeze shorts at will.
\end{definition}

\section{Nickel, March 2022}

\begin{omfigure}
\begin{tikzpicture}
\begin{axis}[width=11cm,height=5.2cm,
  ylabel={thousand USD per tonne},
  tick label style={font=\small},label style={font=\small},
  xmin=0.5,xmax=8.5,ymin=0,ymax=115,grid=major,grid style={gray!25},
  xtick={1,2,3,4,5,6,7,8},
  xticklabels={Fri open,Fri close,Mon open,Mon close,Tue 04:49,Tue 06:08,Tue 07:00,Tue 08:15},
  xticklabel style={rotate=35,anchor=north east,font=\footnotesize},
  table/col sep=comma]
\addplot[omThm,very thick,mark=*,mark size=1.6pt] table[x=i,y=price]{figdata/markets-3/08-metals/nickel.csv};
\node[font=\footnotesize,anchor=east] at (axis cs:5.8,101.4) {peak 101{,}365};
\node[font=\footnotesize,anchor=north west] at (axis cs:4.1,46) {48{,}078};
\node[font=\footnotesize,anchor=south west] at (axis cs:0.8,36) {27{,}080 to 28{,}919};
\node[font=\footnotesize,anchor=north,align=center] at (axis cs:7.5,72) {suspended\\at 08:15};
\end{axis}
\end{tikzpicture}

\omcaption{LME three-month nickel, 4 to 8 March 2022: the prices stated in the Court of Appeal's
judgment, in trading order (the weekend omitted; the levels after 07:00 shown at the judgment's
``consistently above 80{,}000''). All trades after midnight on 8 March were
cancelled. Source: \emph{R (Elliott Associates) v London Metal Exchange} [2024] EWCA Civ
1168.}\label{fig:m3:metals:nickel}
\end{omfigure}

The prices rose in three steps (\cref{fig:m3:metals:nickel}). On Friday 4 March three-month
nickel rose 6.8\%, from 27{,}080 to 28{,}919 dollars, after Russia's invasion of Ukraine a week
earlier; the clearing house called about 2.6 billion dollars of intraday margin, 40\% above its
previous record. On Monday 7 March it rose 66\% close to close, to 48{,}078, a daily move nearly
five times larger than any in the previous twenty years, and margin calls totalled about 7.05
billion dollars; three clearing members failed to pay on time. In the first hours of Tuesday it
rose to 101{,}365. At 04:49, with the price near 60{,}000, the exchange's staff suspended the
price bands because trades were being done above them; the price stayed above 80{,}000 from
07:00 until the suspension.

The court's account of the cause was a short squeeze (One Quant Book 1, chapter 16): large short
positions had been built up by several participants, among them a Chinese nickel producer,
Tsingshan, on the over-the-counter market; as the price rose, shorts had to buy to meet margin,
and others bought in front of them. The market reopened on 16 March, after the producer's banks
had put support in place. The cancellation was challenged in court by trading firms that had
bought during the night; the High Court and then the Court of Appeal (2024) held that the
exchange had acted lawfully.

\begin{dated}{2026-09}{The regulator's finding}\label{dat:m3:metals:fca}
On 19 March 2025 the Financial Conduct Authority fined the London Metal Exchange
\pounds9{,}245{,}900, its first enforcement action against a recognised investment exchange, for
failing to maintain orderly trading in March 2022. It found that only junior staff, not trained
to recognise all causes of disorder, were on duty in the Asian hours, and that they worked to
accommodate the rises, including by disabling price bands. The exchange settled and received a
30\% discount.
\end{dated}

\begin{proposition}[Margin on a squeezed short]\label{prop:m3:metals:margin}
A short of $q$ tonnes marked from price $P_0$ to $P_1$ owes variation margin $q\,(P_1 - P_0)$ in
cash the same day. A short of 10{,}000 tonnes owed 191.59 million dollars on the Monday and
would have owed a further 532.87 million at the Tuesday peak: 724.46 million on a position
whose value had been 289 million dollars at Friday's close.
\end{proposition}

\begin{proof}
$10\,000 \times (48\,078 - 28\,919)$ and $10\,000 \times (101\,365 - 48\,078)$; the position was
worth $10\,000 \times 28\,919$ at Friday's close.
\end{proof}

A producer that has sold its future output forward is economically hedged: the metal it will
produce rises in value with its short. Its margin, though, is paid in cash now, and the metal's
gain arrives only when it is sold. The nickel squeeze is the sharpest example of a hedge that
was sound in value and fatal in liquidity, a theme of \cref{ch:m3:when-markets-break-together}.

\section{Gold: London against New York}

Gold trades in two places with different conventions. In London it trades over the counter,
for settlement in unallocated metal.

\begin{definition}[Loco London, unallocated gold]\label{def:m3:metals:loco}
\emph{Loco London}\index{loco London} is the convention that a price is for gold (or silver)
delivered in London, settled through the London clearing system of unallocated accounts.
\emph{Unallocated gold}\index{unallocated gold} is a balance of gold in an account with a clearing
bank: the holder owns no specific bars, only a claim on the bank for that weight of metal, and
is an unsecured creditor of the bank.
\end{definition}

\begin{definition}[LBMA Gold Price]\label{def:m3:metals:lbma}
The \emph{LBMA Gold Price}\index{LBMA Gold Price} is the benchmark price of spot, unallocated,
\omterm{def:m3:metals:loco}{loco London} gold, set twice each London business day in electronic auctions run by an
independent administrator.
\end{definition}

\begin{dated}{2026-09}{The London auctions}\label{dat:m3:metals:lbma}
The \omterm{def:m3:metals:lbma}{LBMA Gold Price} auctions run at 10:30 and 15:00 London time, administered by ICE Benchmark
Administration, which has run them since March 2015.
\end{dated}

In New York gold trades as futures on COMEX, delivered in 100-ounce bars in approved vaults. The
two markets are linked by the exchange for physical (One Quant Book 1, chapter 21): a trader long
futures and short London gold swaps one for the other, and normally the futures trade at a
premium over London equal to carry and a small fee, about \$1.50 an ounce. In March 2020 the link
broke. Swiss refineries that recast London's 400-ounce bars into the smaller bars New York
delivers closed, and the passenger flights that carry gold were grounded; on 25--26 March April
futures traded more than \$70 above London and then at \$50. The exchange and the London
association responded with a futures contract deliverable in 100-ounce, 400-ounce or one-kilogram
bars.

\begin{omfigure}
\centering
\begin{tikzpicture}[font=\small,
  mk/.style={draw,thick,rounded corners=3pt,align=center,minimum width=3.9cm,minimum height=1.6cm},
  arr/.style={{Stealth[length=2.2mm]}-{Stealth[length=2.2mm]},thick}]
\node[mk,draw=omDef,fill=omDef!8] (ldn) at (0,0) {London OTC\\unallocated, loco London\\400-ounce bars};
\node[mk,draw=omThm,fill=omThm!8] (ny) at (8.4,0) {New York futures\\COMEX vaults\\100-ounce bars};
\draw[arr] (ldn.east) -- node[above=3pt,font=\footnotesize] {exchange for physical} node[below=3pt,font=\footnotesize] {normally about \$1.50} (ny.west);
\node[font=\footnotesize,align=center,anchor=north] at (4.2,-1.2) {physical link: refiners recast bars,\\and passenger flights carry them};
\end{tikzpicture}

\omcaption{The two gold markets and the link between them. The EFP spread stays near its cost
while bars can be recast and flown; in March 2020 both failed at once. Schematic; source of the
normal level and of the 2020 dislocation: Singapore Bullion Market Association,
2020.}\label{fig:m3:metals:gold}
\end{omfigure}

\section{Iron ore}

Iron ore trades on two very different kinds of contract. In Singapore it trades as cash-settled futures and swaps on the monthly average
of a \omterm{def:m3:physical-commodity-markets:pra}{price-reporting agency}'s assessment of 62\% iron-content fines delivered to China; in Dalian
it trades as physically delivered futures in lots of 100 tonnes, open to foreign investors since
May 2018. The spread between them is a market in its own right, with \omterm{def:m3:physical-commodity-markets:basisrisk}{basis risk} of the kind met
in \cref{ch:m3:physical-commodity-markets}.

\section{Tutorial: the prompt calendar and a squeezed short}\label{tut:m3:metals:prompts}

\begin{tutorial}
\textbf{Goal.} Build the prompt-date calendar and compute the margin of a short through the nickel
squeeze. \textbf{End state:} \cref{fig:m3:metals:prompts} and the margins of
\cref{prop:m3:metals:margin}.
\begin{tutsteps}
\item \textbf{\omterm{def:m3:metals:prompt}{Prompt dates}.} Daily to three months, weekly Wednesdays to six months, third
Wednesdays beyond.
\omcode{code/firm/prompts/firm_prompts.py}{59}{79}{The prompt dates of a trade date.}\label{lst:m3:metals:prompts}
\item \textbf{The squeeze.} Variation margin day by day on the judgment's prices.
\omcode{code/markets-3/08-metals/python/m3_metals.py}{28}{34}{Variation margin of a short through 7 and 8 March 2022.}\label{lst:m3:metals:squeeze}
\item \textbf{Run} \texttt{m3\_metals.calendar()}, \texttt{short\_squeeze()} and \texttt{fig\_metals.py}.
\end{tutsteps}
\textbf{What to change next.} Add the exchange's holiday calendar; price a carry from cash to three
months from two quotes and convert it to an annualised implied interest rate.
\end{tutorial}

\section{Build: the prompt calendar and margin estimator}\label{bld:m3:metals:prompts}

\begin{build}
\bfield{Purpose} The miniature firm hedges metal on the LME: it must know which dates exist today,
what a carry costs, and how much cash its positions will need if prices move.
\bfield{Interface} \texttt{prompt\_dates(trade, months, holidays)} returning daily, weekly, monthly,
cash and three-month dates; \texttt{carry(near, far)}; \texttt{variation\_margin(position, old,
new)}; the calendar helpers \texttt{add\_months}, \texttt{lme\_roll}, \texttt{three\_month\_date},
\texttt{third\_wednesday}.
\bfield{Rules} Business days exclude weekends and the given holidays; a \omterm{def:m3:metals:prompt}{prompt date} on a non-business
day moves to the next business day, but a Saturday moves back to the Friday, and a three-month date
pushed into the fourth month falls on the last business day of the third (the LME's Trading
Regulations 8.4.1 and 8.4.2); monthly prompts start after the last weekly one.
\bfield{Acceptance tests} \texttt{code/firm/prompts/tests/}: month arithmetic at month ends; the Saturday,
Sunday and month-end rolls; the
daily, weekly and monthly structure; the sign of a short's margin.
\bfield{Stretch} Load the exchange's published trading calendar; margin under a scenario of
several days' moves with intraday calls; lending-rule limits on a dominant long.
\end{build}

\begin{omsources}
\item Court of Appeal, \emph{R (Elliott Associates) v London Metal Exchange} [2024] EWCA Civ 1168.
\item Financial Conduct Authority, press release of 20 March 2025 and Final Notice to the London
Metal Exchange.
\item London Metal Exchange, prompt date structure; position management and lending rules; warehouse
policy; \emph{Rules and Regulations} (1 June 2026), Part 3, Trading Regulation 8, prompt dates (Internet Archive copy).
\item LBMA, the LBMA Gold Price and loco London; ICE Benchmark Administration.
\item J.~East, ``The day the EFP broke'', Singapore Bullion Market Association, 26 March 2020.
\item SGX iron ore contract rules; Dalian Commodity Exchange, opening of iron ore futures to foreign
investors (2018).
\end{omsources}

\section{Exercises}

\begin{exercise}[$\star$]\label{exo:m3:metals:1}
On a Thursday, which date is tom and which is cash, if no holidays intervene?
\end{exercise}

\begin{exercise}[$\star$]\label{exo:m3:metals:2}
A consumer lends 500 tonnes of copper tom-next. What does it do on each date?
\end{exercise}

\begin{exercise}[$\star$]\label{exo:m3:metals:3}
Why does a surge in \omterm{def:m3:metals:warrant}{cancelled warrants} tend to tighten the nearby dates?
\end{exercise}

\begin{exercise}[$\star\star$]\label{exo:m3:metals:4}
A short of 2{,}000 tonnes of nickel was marked at Friday's close. What variation margin did it owe
on Monday, and what would it have owed at Tuesday's peak?
\end{exercise}

\begin{exercise}[$\star\star$]\label{exo:m3:metals:5}
By how much did nickel rise from Friday's open to Tuesday's peak? Give each day's move.
\end{exercise}

\begin{exercise}[$\star\star$]\label{exo:m3:metals:6}
In March 2020, why could a trader not simply buy London gold at \$1{,}600 and deliver it into
COMEX futures at \$1{,}670?
\end{exercise}

\begin{exercise}[$\star\star\star$]\label{exo:m3:metals:7}
\emph{Coding.} With \texttt{prompt\_dates}, count the daily, weekly and monthly prompts available on
24 September 2026 out to 27 months.
\end{exercise}

\begin{exercise}[$\star\star\star$]\label{exo:m3:metals:8}
\emph{Find the flaw.} ``The producer was fully hedged, so the price spike could not hurt it.''
\end{exercise}

\section{Problem: Tuesday 8 March 2022}

\begin{problem}[{Weekend problem --- a short in the nickel squeeze}]\label{pb:m3:metals:1}
A metals trader is short 10{,}000 tonnes of three-month nickel against expected output. Use the
prices of \cref{fig:m3:metals:nickel}.

\medskip\noindent\textbf{Part I --- The position.}
\begin{enumerate}
\item What was the position worth at Friday's close?
\item What margin did Friday's rise cost?
\item What did Monday's rise cost?
\item What would the Tuesday peak have cost, had the trades stood?
\item Give the total cash drain from Friday's open to the peak.
\end{enumerate}

\medskip\noindent\textbf{Part II --- The market.}
\begin{enumerate}[resume]
\item How large were the clearing house's margin calls on 4 and 7 March?
\item What would members have had to post on 8 March had the trades stood?
\item Why were the price bands suspended at 04:49?
\item Why did the exchange cancel rather than only suspend?
\item Who lost from the cancellation, and what did the courts decide?
\end{enumerate}

\medskip\noindent\textbf{Part III --- The mechanism.}
\begin{enumerate}[resume]
\item Why does a known large short invite others to buy?
\item How does a lending rule limit squeezes on the exchange, and why did it not stop this one?
\item Why did the squeeze start on the over-the-counter market?
\item What role did thin Asian-hours trading play?
\item What made the market reopen on 16 March?
\end{enumerate}

\medskip\noindent\textbf{Part IV --- Judgement.}
\begin{enumerate}[resume]
\item What should the trader have held besides the short to survive the week?
\item What would the regulator's findings change in the exchange's operations?
\item Should an exchange ever cancel trades that were not errors?
\item State the \emph{named result}: the margin a 10{,}000-tonne short owed from Monday's close to
the peak, and in total from Friday's close.
\item In one sentence: what kills a hedged producer in a squeeze?
\end{enumerate}
\end{problem}

\section{Interview questions}

\begin{interviewq}[$\star$ \iqroles{trader}]\label{iq:m3:metals:1}
How is an LME contract different from a CME monthly future?
\end{interviewq}

\begin{interviewq}[$\star$ \iqroles{trader, risk}]\label{iq:m3:metals:2}
What is a short squeeze, and what did the nickel episode add to your understanding of one?
\end{interviewq}

\begin{interviewq}[$\star\star$ \iqroles{researcher}]\label{iq:m3:metals:3}
Stocks on warrant fall by half in a month. What happens to the curve, and how would you trade it?
\end{interviewq}

\begin{interviewq}[$\star\star$ \iqroles{risk}]\label{iq:m3:metals:4}
Design the liquidity stress test for a producer hedging its output with exchange futures.
\end{interviewq}

\begin{interviewq}[$\star\star$ \iqroles{trader}]\label{iq:m3:metals:5}
Why did the gold EFP blow out in March 2020, and what closed it?
\end{interviewq}

\begin{interviewq}[$\star\star\star$ \iqroles{developer, risk}]\label{iq:m3:metals:6}
You run an exchange's surveillance. What signals would have warned you on the night of 7--8 March
2022, and what would you have automated?
\end{interviewq}
